%% A crossing may name the semantic wire generated for an unmatched typed
%% port.  Every generated leg enters the saved-route crossing pipeline.
\documentclass{standalone}
\usepackage{tenkz}
\makeatletter
\ExplSyntaxOn
\file_input:n { tenkz-kernel.code.tex }
\int_new:N \g__tenkz_test_promoted_port_open_int
\cs_new_eq:NN
  \__tenkz_test_port_open_stroke:nn
  \__tenkz_kernel_r_wire_stroke:nn
\cs_set_protected:Npn \__tenkz_kernel_r_wire_stroke:nn #1#2
  {
    \__tenkz_model_get:nnN
      { \l__tenkz_kernel_r_wire_record_tl } {origin} \l_tmpa_tl
    \str_if_eq:VnT \l_tmpa_tl {port-open}
      { \int_gincr:N \g__tenkz_test_promoted_port_open_int }
    \__tenkz_test_port_open_stroke:nn {#1} {#2}
  }
\AtEndDocument
  {
    \int_compare:nNnF { \g__tenkz_test_promoted_port_open_int } = {1}
      {
        \errmessage
          {crossing~police~did~not~register~exactly~one~open~port~route}
      }
  }
\ExplSyntaxOff
\makeatother
\begin{document}
\tenkzkernel
\begin{tenkz}[rows={wire}, cols=3, bonds=none]
  \tn[at=(1,2), name=P, ports={90:physical}]{P}
  \tnwire[
    kind=string,
    name=probe,
    via={1.5 n of (1,2)},
    cross={under at crossing of probe and port-open-1}
  ]{1.5 n of (1,1)}{1.5 n of (1,3)}
  \tnwire[
    kind=string,
    name=near-probe,
    via={0.8 n of (1,2)},
    cross={under at crossing of near-probe and port-open-1}
  ]{0.8 n of (1,1)}{0.8 n of (1,3)}
\end{tenkz}
\end{document}
